% !TeX root = ../main.tex
\section{Reference response and limiting cases}
\label{sec:compatibility}

\subsection{Reference stress and solid storage}
At the undeformed zero-pressure state, the logarithmic derivative reduces
to the identity on symmetric tensors. Using the drained stiffness relation
in \eqref{eq:anisotropic-biot-explicit} gives
\begin{equation}
 \mathbf B_0=\mathbf I-\mathbb{C}^d:\mathbb{C}_s^{-1}:\mathbf I.
 \label{eq:reference-biot-compatibility}
\end{equation}
This is the anisotropic reference relation associated with homogeneous
mineral compression. Its isotropic form is
\(\mathbf B_0=(1-K/K_s)\mathbf I\).
The change in pore volume with pressure at fixed skeleton deformation
also follows from the mineral equation. Define the reference solid
contribution to storage by
\begin{align}
 S_s&=-\left.\phi_{s0}\pder{\bar J}p\right|_{\mathbf F=\mathbf I,p=0},
 \label{eq:reference-solid-storage-definition}\\
 &=\frac{\phi_{s0}}{K_s}\left(1-\frac K{\phi_{s0}K_s}\right),
 \label{eq:reference-solid-storage}\\
 &=\phi_{s0}\mathbf I:\mathbb{C}_s^{-1}:\mathbf I
 -\mathbf I:\mathbb{C}_s^{-1}:\mathbb{C}^d:\mathbb{C}_s^{-1}:\mathbf I.
 \label{eq:reference-storage-compatibility}
\end{align}
Fluid compressibility supplies a separate fluid storage contribution.
The stress and solid storage relations here come from the same energy and
the same mineral-volume response.

\subsection{Isotropic mineral}
For an isotropic mineral,
\(\mathbb{C}_s:\mathbf I=3K_s\mathbf I\). The mineral volume--shape
term vanishes and \eqref{eq:anisotropic-mineral-eos} becomes
\begin{equation}
 K_s\ln\bar J+\left(1-\frac K{\phi_{s0}K_s}\right)p\bar J
                   -\frac K{\phi_{s0}}\ln J=0.
 \label{eq:reconstructed-isotropic-source-eos}
\end{equation}
The Biot tensor is spherical,
\begin{equation}
 \mathbf B=\left[1-\frac{K\bar J}
 {J\left[K_s+\left(1-K/(\phi_{s0}K_s)\right)p\bar J\right]}\right]
 \mathbf I.
 \label{eq:reconstructed-isotropic-source-biot}
\end{equation}
These are the elastic scalar relations in the companion manuscript.
Eliminating the two constituent volume factors in Gajo's logarithmic model
also gives \eqref{eq:reconstructed-isotropic-source-eos}
\citep[eqs.~(3.27), (3.32), and (3.34)]{gajo2010}. The correspondence concerns
its volumetric law. Here the full Hooke law and volume-only conformal distention
require the drained shear stiffness to be \(\phi_{s0}\) times the mineral
shear stiffness, as follows from \eqref{eq:drained-stiffness-restriction}.

\subsection{Equal confining pressure and pore pressure}
An unjacketed specimen has confining pressure equal to pore pressure. The
solid and mixture should then have the same homogeneous stretch; their
representations may differ by the internal rotation.
Choose a state satisfying
\begin{equation}
 a=1,\qquad J=\bar J,\qquad
 \mathbb{C}_s:\mathbf\varepsilon=-pJ\mathbf I.
 \label{eq:reconstructed-finite-unjacketed-path}
\end{equation}
The logarithmic strain can be anisotropic because \(\mathbb{C}_s^{-1}:\mathbf I\)
need not be spherical. The logarithmic stress transformation maps
\(\mathbf I\) to \(\mathbf I\), so the mineral Cauchy stress is
\(-p\mathbf I\). Distention stress vanishes, the mineral-volume equation
is satisfied, and the phase balance gives
\begin{equation}
 \mathbf\sigma=-p\mathbf I,\qquad \phi_s=\phi_{s0}.
 \label{eq:finite-unjacketed-stress-fraction}
\end{equation}
These identities hold on the selected elastic branch. They verify finite
homogeneous mineral compression, in addition to the reference linear
relations.
